\clearpage
\fancyhead[L]{\small KOMET.}
\hypertarget{native-464-comet}{}
\begin{center}{\large KOMET.}\end{center}

\begin{center}
\rule{0.55\linewidth}{0.4pt}\\[2pt]
Astronomische Nachrichten Band VI. Nr. 123. Seite 43\,..\,44. 1827 September.\\[2pt]
\rule{0.55\linewidth}{0.4pt}
\end{center}

\begin{center}Beobachtungen des von \textsc{Pons} im Luchs entdeckten Kometen.\end{center}

Hofrat \textsc{Gauss} hat mir die folgenden seiner Meridianbeobachtungen geschickt.

\begin{center}
\begin{tabular}{l|l|l|l}
 & & Rektasz. Kom. & $\delta$ Kom. \\
\hline
1827 Aug. 20 & $10^{\rm u}28^{\rm m}11^{\rm s}1$ & $125^\circ29'41''9$ & $+56^\circ21'56''1$ \\
\hfill 21 & $37\;53.0$ & $128\;54\;42.5$ & $54\;58\;38.2$ \\
\hfill 22 & $47\;27.0$ & $132\;17\;54.6$ & $53\;23\;1.1$ \\
\end{tabular}
\end{center}

Aus diesen Beobachtungen hat Herr \textsc{Peters} folgende Bahn berechnet:

\Needspace{9\baselineskip}
\SourceLeaf{465}
Zeit des Perihels 1827 September $12.2632$ Altona.

\begin{center}
\begin{tabular}{l l l}
$P$ & \mbox{.\,.\,.\,.\,.\,.\,.\,.\,.\,.\,.\,.\,.\,.\,.\,.\,.\,.\,.} & $255^\circ25'50''$ \\
$\AscendingNode$ & \mbox{.\,.\,.\,.\,.\,.\,.\,.\,.\,.\,.\,.\,.\,.\,.\,.\,.\,.\,.} & $150\;11\;40$ \\
$i$ & \mbox{.\,.\,.\,.\,.\,.\,.\,.\,.\,.\,.\,.\,.\,.\,.\,.\,.\,.\,.} & $54\;27\;5\mathrm{c}$ \\
$\log q$ & \mbox{.\,.\,.\,.\,.\,.\,.\,.\,.\,.\,.\,.\,.\,.\,.\,.\,.\,.} & $9.21395$ \\
\end{tabular}
\end{center}

Rückläufig.

Diese Elemente geben die mittlere Beobachtung mit Abweichungen von $+5''$ in der Länge und $-4''$ in der Breite wieder.
\begin{flushright}\textsc{Schumacher.}\end{flushright}

